
\label{sec:compositional_generalization}

In the next set of experiments, we study how well \MethodName{} can perform new tasks by compositionally generalizing over the skills seen in training. This has been viewed as a kind of ``grand challenge'' for robotic foundation models: while prior models have demonstrated generalization over semantic concepts (e.g., reaching for an object with an unseen textual label), performing new tasks has proven elusive.

We found that for some short-horizon tasks, \MethodName{} can work well completely out of the box, despite no data being collected explicitly for these tasks. These tasks involve manipulating unfamiliar objects in new ways, such as wiping down a pair of headphones with a cloth, or spinning a desk fan (Fig.~\ref{fig:task_generalization}). For longer-horizon, more complex tasks, we find that we can actually use the instruction following capabilities of \MethodName{} to ``coach'' it through the task with language. This provides an exciting new way to teach \MethodName{} new tasks without collecting any additional training data.

\noindent \textbf{\MethodName{} can perform new short-horizon tasks out of the box. }
We find that \MethodName{} can perform short-horizon tasks such as pressing the plunger on a french press, scooping rice into a rice cooker, wiping down common office objects, and spinning various articulated items directly out of the box, despite no robot data being collected specifically for any of these tasks (Fig.~\ref{fig:task_generalization}). This is quite exciting --- since \MethodName{} can flexibly compose a wide variety of skills, we can often get it to perform new, simple tasks by just prompting it to get the desired behaviors.

\noindent \textbf{\MethodName{} can be coached to perform new longer horizon tasks purely with language. }
While \MethodName{} can be directly prompted to perform new short-horizon tasks, simply asking the robot to perform an unseen long-horizon task like cooking a sweet potato would not work: although \MethodName{} exhibits a high level of out-of-the-box generalization, these tasks are simply too complex, requiring up to 5 minutes of interaction with multiple stages. However, \MethodName{}'s language following abilities provide us with an exciting new path to teach the model such tasks: instead of providing demonstration data for each complex skill that we want the robot to learn, we can instead ``coach'' the robot to perform the new task with language, a bit like how one might teach the task to a person (Fig.~\ref{fig:air_fryer_coaching}). We set up several realistic multi-stage kitchen tasks: (1) ``Loading an Air Fryer'': using an air fryer to cook a sweet potato; (2) ``Unloading an Air Fryer'': dumping items out of an air fryer, and (3) ``Toasting a Bagel'': using a toaster to toast up a bagel. In each case, our robot data did not contain any training episodes for this task, although similar appliances were seen in different contexts in human data and external datasets. However, a person can walk the robot through task with step-by-step instructions, such as ``pick up the sweet potato'' and ``open the air fryer.'' We present the results of coaching \MethodName{} and comparisons in Fig.~\ref{fig:long_horizon_task_generalization}. Critically, none of the models have any \emph{action-level} data of these particular tasks, and in the ``coaching'' episodes, the environment and task are entirely unseen. We find that \MethodName{} can be coached much more effectively than prior methods to perform all of these tasks, and can perform even more effectively when conditioned on generated subgoal images.

\noindent \textbf{Coaching data can endow \MethodName{} with new capabilities. }
Since \MethodName{} can be coached to perform new tasks, we can actually use the step-by-step instructions in the coaching data to train a high-level language policy which can prompt \MethodName{} with the appropriate language instructions as it performs the task. This gives \MethodName{} the ability to perform fully unseen, long-horizon tasks without collecting any additional teleoperation data. We show the results of this experiment in Fig.~\ref{fig:coaching} for five different tasks. For all of these tasks, we find that we can successfully train an autonomous policy (\MethodName{} (autonomous)) that can roughly match the performance of the coaching episodes (\MethodName{} (coaching)) that were collected by prompting the model to perform the task.



\subsection{Can \texorpdfstring{\MethodName{}}{pi0.7} learn effectively from diverse and mixed-quality data?}
\label{sec:scaling}
In our last set of experiments, we perform a set of controlled ablation studies to understand whether \MethodName{} can effectively leverage large and diverse datasets, and whether its performance improves with dataset diversity. These questions are difficult to answer definitively, since the performance of such a model depends on a large number of factors, and very large datasets are very difficult to slice cleanly so as to ablate ``diversity". To provide some understanding of these questions, we first study whether \MethodName{} continues to improve on seen tasks when trained on increasingly larger but more mixed-quality datasets. Then, we study whether \MethodName{} can leverage datasets with high task diversity to drive improvements in generalization.

\noindent \textbf{\MethodName{} can effectively learn from mixed-quality data: }
Effectively learning from diverse robot data has so far been a significant challenge in training robotic policies. Designers will often carefully filter or curate the data to pull high-quality datasets with consistent strategies \cite{hejna2025robot, li2025gr}. However, data filtering is laborious, task-specific, and ends up throwing away a lot of valuable information. In these experiments, we aim to answer: can \MethodName{} learn more from data with diverse manipulation strategies?

To study this question, we consider the ``Laundry (T-Shirts and Shorts)'' task that we studied in Fig.~\ref{fig:distillation_results}. We annotated the data collected by human operators based on fold quality and speed, and split our dataset into 4 buckets consisting of (1) the top 30\% by quality and speed, (2) the top 50\%, (3) the top 80\%, and (4) all of the data. We then train new \MethodName{} models from scratch for data from each of the four buckets, using episode metadata or without it (8 models in total).
We find that while \MethodName{} (without metadata) can actually get worse when trained on larger, mixed-quality datasets, \MethodName{} (with metadata) is able to continuously improve as we train on more data, even though the dataset size increases corresponds to a decrease in average data quality (Fig.~\ref{fig:generalization_and_conditioning}, left). This suggests that our diverse prompting method effectively makes the model design \emph{more scalable}, in the sense that it benefits more from larger datasets, even when these larger datasets actually consist of lower quality data that harm models trained in the usual way. Episode metadata effectively disambiguates the different data quality and strategies within the large-scale datasets during \MethodName{} training, and enables prompting for the desired mode of behavior at test time.

\noindent \textbf{Can \MethodName{} translate increased dataset diversity into better generalization performance?}

To study this question, we compare \MethodName{} with the following ablations on some of the short-horizon unseen tasks from Fig.~\ref{fig:task_generalization}.
\begin{itemize}
    \item \MethodName{} (w/o most diverse 20\%): \MethodName{} but with the 20\% of our data with the highest task diversity removed.
    \item \MethodName{} (w/o random 20\%): \MethodName{} with a randomly sampled 20\% of our data removed to serve as a data-controlled comparison to \MethodName{} (w/o most diverse 20\%)
\end{itemize}

The comparison between \MethodName{} (w/o most diverse 20\%) and \MethodName{} (w/o random 20\%) allows us to understand the impact of data with high task diversity in a controlled way, as both models are trained on the same quantity of data. We find that across all tasks, \MethodName{} and \MethodName{} (w/o random 20\%) significantly outperform \MethodName{} (w/o most diverse 20\%), demonstrating that \MethodName{} is able to effectively ingest data with high task diversity and translate that data into performance improvements on short-horizon, unseen tasks (Fig.~\ref{fig:generalization_and_conditioning}, right).

